% Walter Thomas Moya Araya — CV final, 8 de octubre de 2026.
% Desde la raíz del repositorio: tectonic --outdir public/cv documents/cv/walter_thomas_moya_araya_cv_es.tex
% También compatible con pdfLaTeX y XeLaTeX (paquetes estándar de TeX Live).
% Una columna; contacto en el cuerpo; texto Unicode; sin imágenes ni tablas.
\documentclass[10pt,a4paper]{article}
\usepackage{iftex}
\ifPDFTeX
  \usepackage[T1]{fontenc}
  \usepackage[utf8]{inputenc}
  \input{glyphtounicode}
  \pdfgentounicode=1
  \usepackage{lmodern}
  \renewcommand{\familydefault}{\sfdefault}
\else
  \usepackage{fontspec}
  \setmainfont{lmsans10-regular.otf}[
    BoldFont=lmsans10-bold.otf,
    ItalicFont=lmsans10-oblique.otf,
    BoldItalicFont=lmsans10-boldoblique.otf,
    Ligatures=NoCommon]
\fi
\usepackage[spanish,es-nodecimaldot]{babel}
\usepackage[left=1.4cm,right=1.4cm,top=1.2cm,bottom=1.2cm]{geometry}
\usepackage{xcolor}
\usepackage{enumitem}
\usepackage{titlesec}
\usepackage{microtype}
\usepackage{ragged2e}
\usepackage{hyperref}
\definecolor{navy}{HTML}{193D53}
\definecolor{muted}{HTML}{465462}
\hypersetup{
  colorlinks=true,urlcolor=navy,linkcolor=navy,
  pdftitle={Walter Thomas Moya Araya | Ciencia de Datos y Análisis Cuantitativo},
  pdfauthor={Walter Thomas Moya Araya},
  pdfsubject={CV | Data Scientist Junior, Data Analyst Junior, Análisis Cuantitativo},
  pdfkeywords={Python, SQL, Power BI, BigQuery, Machine Learning, K-means, ETL},
  pdflang={es-CL}}
\urlstyle{same}
\pagestyle{empty}
\setlength{\parindent}{0pt}
\setlength{\parskip}{1.6pt}
\setlength{\emergencystretch}{1.5em}
\setlist[itemize]{leftmargin=1.1em,label=\textbullet,itemsep=1.3pt,topsep=2pt,parsep=0pt,partopsep=0pt}
\titleformat{\section}{\fontsize{10.7}{12}\selectfont\bfseries\color{navy}}{}{0pt}{}[\vspace{-0.5pt}\color{navy!25}\titlerule]
\titlespacing*{\section}{0pt}{8pt}{4.5pt}
\newcommand{\entry}[2]{\par\addvspace{3pt}\textbf{#1}\hfill{\fontsize{9.5}{11}\selectfont\color{muted}#2}\par\nopagebreak}
\newcommand{\project}[3]{\par\addvspace{3pt}\textbf{\href{#1}{#2}}\hfill{\fontsize{9.5}{11}\selectfont\color{muted}#3}\par\nopagebreak}

\begin{document}
\fontsize{10.7}{12.5}\selectfont
\RaggedRight
\hyphenpenalty=10000
\exhyphenpenalty=10000

{\fontsize{23}{26}\selectfont\bfseries\color{navy}Walter Thomas Moya Araya}\par
{\fontsize{11}{13}\selectfont\bfseries Data Scientist Junior | Data Analyst Junior | Análisis Cuantitativo}\par
{\fontsize{9.7}{11.5}\selectfont
Santiago, Chile \enspace|\enspace \href{tel:+56933669343}{+56 9 3366 9343} \enspace|\enspace
\href{mailto:cg.walter.ma@gmail.com}{cg.walter.ma@gmail.com}\par
\href{https://waltermoya.vercel.app/}{waltermoya.vercel.app} \enspace|\enspace
\href{https://github.com/walterm2482}{github.com/walterm2482} \enspace|\enspace
\href{https://www.linkedin.com/in/walter-moya-araya-a211b9307/}{linkedin.com/in/walter-moya-araya-a211b9307}}

\section{Perfil profesional}
Ingeniero Civil Industrial y Magíster en Ciencias de la Ingeniería, con experiencia en ETL y calidad de datos
en PwC Chile y desarrollo cuantitativo remoto en ClickAlgo. Combino Python, SQL y estadística para construir
pipelines reproducibles, evaluar modelos y transformar datos en herramientas de decisión.

\section{Habilidades técnicas}
\textbf{Datos y BI:} Python (Pandas, NumPy), SQL, R, Excel, Power BI, ETL y análisis exploratorio de datos.\par
\textbf{Machine learning:} scikit-learn, TensorFlow/Keras, clustering, PCA, series de tiempo y validación temporal.\par
\textbf{Cloud y desarrollo:} Google Cloud (BigQuery, Data Fusion, Cloud Storage), C\#/.NET, Git, Docker y MLflow.

\section{Proyectos destacados}
\project{https://clickalgo.com/k-means}{K-Means Indicator v2 | Machine learning aplicado}{ClickAlgo · 2025}
Desarrollé un indicador para identificar regímenes de mercado mediante clustering K-means, PCA y variables
de precio, volatilidad y momentum. Refactoricé el código e incorporé visualizaciones y alertas por Telegram.

\project{https://ctrader.com/products/438}{FSG Ultimate / Smart Grid Ultimate v2.1 | Automatización y riesgo}{cTrader Store · 2024--2026}
Desarrollé y publiqué un robot de reversión a la media en C\#/.NET, con ejecución sobre velas cerradas,
objetivos según volatilidad (ATR) y tres capas de control de riesgo para gestionar exposición y pérdidas.

\project{https://clickalgo.com/smart-portfolio-architect}{Smart Portfolio Architect | Diversificación de portafolios}{ClickAlgo · 2024}
Automaticé la selección de activos por mínima correlación en C\# para cTrader. La herramienta analiza datos
históricos y compara combinaciones de instrumentos para apoyar decisiones de diversificación.

\project{https://github.com/walterm2482/portfolio-optimizer}{Portfolio Optimizer | Python, Dash y backtesting}{GitHub · 2025}
Construí una aplicación con cuatro métodos de asignación de activos, backtesting con rebalanceo y
visualizaciones de rentabilidad, drawdown y correlación para comparar estrategias.

\section{Experiencia profesional}
\entry{Desarrollador cuantitativo y de IA | Independiente}{10/2024 -- actualidad}
\begin{itemize}
  \item Construí una plataforma de investigación en Python con datos de mercado, ingeniería de variables,
  modelos predictivos y backtesting; automaticé controles de calidad y trazabilidad de resultados.
  \item Apliqué validación temporal y evaluación fuera de muestra para reducir fuga de información entre
  entrenamiento y evaluación de modelos de series de tiempo.
\end{itemize}

\entry{Programador junior (Junior Programmer) | ClickAlgo, Reino Unido · Remoto}{11/2023 -- 10/2024}
\begin{itemize}
  \item Desarrollé indicadores cuantitativos en C\#/.NET con clustering, filtros de Kalman y análisis de volatilidad;
  entregué módulos probados y documentados y brindé soporte a usuarios internacionales.
  \item Publiqué un artículo técnico sobre portafolios de mínima correlación con Pearson, Spearman y Kendall.
\end{itemize}

\entry{Analista en formación (Analyst Trainee) | PwC Chile}{08/2022 -- 11/2022}
\begin{itemize}
  \item Optimicé cargas históricas de inventario de cientos de miles de registros a \textbf{menos de 30 segundos por archivo}.
  \item Apoyé pipelines ETL en Google Cloud con Python y APIs REST, validación de esquemas y controles
  de calidad para alimentar dashboards en Power BI.
\end{itemize}

\section{Educación}
\entry{Magíster en Ciencias de la Ingeniería, mención Industrias}{08/2021 -- 06/2023}
Universidad Diego Portales · \textbf{Máxima distinción} · Tesis: 6,8/7,0.\par
Tesis: detección y seguimiento de maquinaria de construcción con YOLOv8 y BoT-SORT.\par
\entry{Ingeniería Civil Industrial}{03/2016 -- 07/2023}
Universidad Diego Portales · Titulado con distinción.

\section{Certificaciones seleccionadas}
MLOps Specialization — Duke University (2026) · Machine Learning — DeepLearning.AI / Stanford (2025).\par
IBM Data Science Professional Certificate (2025) · Google Data Analytics Professional Certificate (2025).
\end{document}
